%% Dados do exercício ia-6.1 — esqueleto de dissertação
%% Matriz curricular informada: Ciência de Dados e Inteligência Artificial.
%% A matriz curricular não foi presumida como área ou linha de pesquisa.

%% Identificação
\titulo{Exposição ocupacional à inteligência artificial e transições no mercado de trabalho brasileiro}
\subtitulo{evidências de painel da PNAD Contínua após o lançamento público do ChatGPT}
\titulotraduzido{Occupational exposure to artificial intelligence and labor market transitions in Brazil}

\autor{Ricardo Carvalho da Costa}
\autorficha{COSTA, Ricardo Carvalho da}

\orientador{Prof. Dr. Danny de Castro Soares}
\coorientador{}

%% Vínculo institucional
\curso{Mestrado Profissional em Administração Pública}
\programa{Administração Pública}
\area{}
\linhapesquisa{}

\cidade{Brasília}
\local{Brasília-DF}
\data{2026}

%% Defesa — informações ainda não confirmadas
\datadefesa{a definir}

\membrobanca{Prof. Dr. Danny de Castro Soares}{Orientador}
\membrobanca{a definir}{Examinador interno}
\membrobanca{a definir}{Examinador externo}

%% Ficha catalográfica — pendente de elaboração oficial
\cutter{}
\numeropaginas{a definir}
\cdd{a definir}
\assuntos{Inteligência artificial. Mercado de trabalho. Mobilidade ocupacional. Informalidade. PNAD Contínua.}

%% Palavras-chave
\palavraschave{Inteligência artificial. Mercado de trabalho. Mobilidade ocupacional. Informalidade. PNAD Contínua.}
\keywords{Artificial intelligence. Labor market. Occupational mobility. Informality. Brazilian household survey.}